\documentclass[12pt,a4paper]{report}

% ============================================================
% PACKAGES
% ============================================================
\usepackage[utf8]{inputenc}
\usepackage[T1]{fontenc}
\usepackage[french]{babel}
\usepackage{geometry}
\usepackage{graphicx}
\usepackage{xcolor}
\usepackage{titlesec}
\usepackage{fancyhdr}
\usepackage{hyperref}
\usepackage{enumitem}
\usepackage{tabularx}
\usepackage{booktabs}
\usepackage{array}
\usepackage{longtable}
\usepackage{listings}
\usepackage{caption}
\usepackage{setspace}
\usepackage{tocloft}
\usepackage{float}
\usepackage{multirow}
\usepackage{tikz}
\usepackage{pifont}

% ============================================================
% CONFIGURATION
% ============================================================
\geometry{left=2.5cm, right=2.5cm, top=2.5cm, bottom=2.5cm}
\setstretch{1.3}

% Colors
\definecolor{attijariOrange}{HTML}{F7941D}
\definecolor{attijariDark}{HTML}{1A1A1E}
\definecolor{attijariGray}{HTML}{6B6C78}
\definecolor{codeBackground}{HTML}{F5F6F8}
\definecolor{greenCheck}{HTML}{16A34A}
\definecolor{blueProgress}{HTML}{2563EB}

% Hyperlinks
\hypersetup{
    colorlinks=true,
    linkcolor=attijariDark,
    urlcolor=attijariOrange,
    citecolor=attijariDark,
}

% Section formatting
\titleformat{\chapter}[display]
    {\normalfont\Large\bfseries\color{attijariDark}}
    {\chaptertitlename\ \thechapter}{14pt}{\LARGE}
\titleformat{\section}
    {\normalfont\large\bfseries\color{attijariDark}}
    {\thesection}{1em}{}
\titleformat{\subsection}
    {\normalfont\normalsize\bfseries\color{attijariDark}}
    {\thesubsection}{1em}{}

% Header/Footer
\pagestyle{fancy}
\fancyhf{}
\fancyhead[L]{\small\textcolor{attijariGray}{Rapport de Mi-Parcours --- PFE 2026}}
\fancyhead[R]{\small\textcolor{attijariGray}{Attijari bank}}
\fancyfoot[C]{\thepage}
\renewcommand{\headrulewidth}{0.4pt}

% Code listing style
\lstset{
    backgroundcolor=\color{codeBackground},
    basicstyle=\ttfamily\small,
    breaklines=true,
    frame=single,
    rulecolor=\color{attijariGray!30},
    tabsize=2,
    captionpos=b,
}

% Custom checkmark and progress symbols
\newcommand{\done}{\textcolor{greenCheck}{\ding{51}}}
\newcommand{\inprog}{\textcolor{blueProgress}{$\circlearrowright$}}

% ============================================================
% DOCUMENT START
% ============================================================
\begin{document}

% ============================================================
% TITLE PAGE
% ============================================================
\begin{titlepage}
    \centering
    \vspace*{1cm}
    
    {\large\textbf{République Tunisienne}}\\
    {\large\textbf{Ministère de l'Enseignement Supérieur}}\\[0.3cm]
    {\large\textbf{et de la Recherche Scientifique}}\\[1.5cm]
    
    \rule{\textwidth}{1.5pt}\\[0.4cm]
    {\Huge\bfseries\textcolor{attijariDark}{Rapport de Mi-Parcours}}\\[0.2cm]
    {\Large\bfseries\textcolor{attijariDark}{Projet de Fin d'Études}}\\[0.3cm]
    \rule{\textwidth}{1.5pt}\\[1cm]
    
    {\LARGE\bfseries\textcolor{attijariOrange}{Système de Détection d'Opportunités}}\\[0.2cm]
    {\LARGE\bfseries\textcolor{attijariOrange}{d'Amélioration Continue des Processus}}\\[0.2cm]
    {\LARGE\bfseries\textcolor{attijariOrange}{Bancaires via l'Intelligence Artificielle}}\\[0.2cm]
    {\LARGE\bfseries\textcolor{attijariOrange}{et la RPA}}\\[1.5cm]
    
    \begin{tabular}{ll}
        \textbf{Réalisé par :} & Meriam \\[0.2cm]
        \textbf{Organisme d'accueil :} & Attijari bank Tunisie \\[0.2cm]
        \textbf{Encadrant Entreprise :} & Firas Elabed (Attijari bank) \\[0.2cm]
        \textbf{Encadrant SESAME :} & Wajih Ben Belgacem \\[0.2cm]
        \textbf{Sujet :} & Sujet 21 \\[0.2cm]
        \textbf{Spécialité :} & Génie Logiciel \\[0.2cm]
        \textbf{Durée :} & 8 semaines (23 Mars -- 17 Mai 2026) \\[0.2cm]
        \textbf{Équipe :} & Rationalisation, Reengineering \\
        & \& Efficacité Opérationnelle \\
    \end{tabular}
    
    \vfill
    
    {\large\textcolor{attijariGray}{Année Universitaire 2025 -- 2026}}
    
\end{titlepage}

% ============================================================
% TABLE OF CONTENTS
% ============================================================
\tableofcontents
\thispagestyle{empty}
\newpage

% ============================================================
% CHAPTER 1: INTRODUCTION
% ============================================================
\chapter{Introduction}

Le secteur bancaire tunisien connaît une transformation numérique profonde, portée par des exigences croissantes en matière d'efficacité opérationnelle, de qualité de service et de conformité réglementaire. Dans ce contexte hautement concurrentiel, les établissements bancaires cherchent à optimiser leurs processus internes, à réduire les délais de traitement des incidents et à améliorer continuellement la satisfaction de leurs clients et collaborateurs.

Le présent rapport constitue le bilan de mi-parcours du projet de fin d'études réalisé au sein de l'équipe \textbf{Rationalisation, Reengineering et Efficacité Opérationnelle} d'Attijari bank Tunisie, filiale du groupe Attijariwafa bank --- leader bancaire panafricain présent dans plus de 25 pays.

Ce projet vise à concevoir et développer un \textbf{système intelligent de détection et d'amélioration continue des processus bancaires}, exploitant les technologies d'intelligence artificielle (NLP, LSTM, KNN) et d'automatisation robotisée des processus (RPA).

Attijari bank enregistre quotidiennement des centaines de réclamations issues de ses clients et collaborateurs internes. Ces réclamations documentent des problèmes variés : erreurs lors du traitement de virements SWIFT, timeouts sur les plateformes de paiement, anomalies dans les systèmes d'information, ou encore dysfonctionnements dans les processus de gestion des comptes. L'historique cumulé de ces réclamations et de leurs actions correctives constitue un corpus de données considérable, mais demeurant largement sous-exploité.

\vspace{0.5cm}
\noindent\fbox{\parbox{\dimexpr\textwidth-2\fboxsep-2\fboxrule}{
\textbf{Problématique centrale :} Comment exploiter intelligemment l'historique des réclamations bancaires et des actions correctives associées afin de détecter automatiquement les anomalies dans les processus, de prédire les incidents avant qu'ils n'impactent les clients, de recommander les actions correctives les plus adaptées, et d'en automatiser l'exécution --- le tout accessible via une interface web sécurisée pour les équipes opérationnelles d'Attijari bank~?
}}
\vspace{0.5cm}

La solution développée repose sur une architecture complète couvrant l'ensemble de la chaîne de valeur~: un backend Python intégrant des techniques avancées de traitement du langage naturel (NLP), d'apprentissage profond (LSTM) et d'automatisation robotisée des processus (RPA), et un frontend développé en Angular 21 et TypeScript offrant aux équipes opérationnelles une interface web intuitive, sécurisée et accessible.

Ce rapport rend compte de l'avancement des travaux réalisés Durant les \textbf{quatre premières semaines} du stage (23 mars -- 17 avril 2026), soit la moitié de la durée totale du projet.

% ============================================================
% CHAPTER 2: TRAVAUX RÉALISÉS — AXE DÉVELOPPEMENT
% ============================================================
\chapter{Travaux Réalisés --- Axe Développement}

L'axe développement constitue le pilier technique sur lequel repose l'ensemble de la solution. Les travaux réalisés au cours des quatre premières semaines couvrent la conception de l'architecture frontend, la mise en place de la base de données, l'implémentation du système d'authentification, et le développement du tableau de bord opérationnel.

\section{Architecture et Technologies Adoptées}

L'architecture de développement repose sur une approche multi-couche robuste et industrielle, découplée et performante~:

\begin{table}[H]
\centering
\caption{Stack technologique retenue}
\begin{tabularx}{\textwidth}{|l|l|X|}
\hline
\textbf{Couche} & \textbf{Technologie} & \textbf{Justification} \\
\hline
Frontend (Production) & Angular 21 · TypeScript & Architecture SPA modulaire réactive, routing avancé et types statiques \\
\hline
Backend (Production) & Python · FastAPI & Framework asynchrone haute performance, intégration directe des modèles IA (LSTM, KNN) \\
\hline
Base de données & PostgreSQL & Support transactionnel complet, contraintes de clés étrangères et intégrité référentielle \\
\hline
Design System & Inter · CSS Variables & Identité visuelle premium Attijari bank, palette de couleurs cohérente \\
\hline
Automatisation & UiPath RPA & Robots logiciels pour l'exécution automatique des workflows de remédiation \\
\hline
Conteneurisation & Docker & Déploiement unifié et portabilité garantie sur n'importe quel environnement client \\
\hline
\end{tabularx}
\end{table}

\section{Développement de l'Interface Frontend (Angular 21)}

L'interface utilisateur a été conçue et développée en Single Page Application (SPA) sous Angular 21, offrant une expérience utilisateur fluide, réactive et hautement sécurisée pour les équipes opérationnelles.

\subsection{Page d'Accueil}

La page d'accueil constitue le point d'entrée de la plateforme. Elle comprend~:

\begin{itemize}[noitemsep]
    \item \textbf{Section Hero} avec animations dynamiques (orbes, grille de fond, gradient text), statistiques clés (70\% d'automatisation, disponibilité 24/7, -45\% de temps de résolution des anomalies)
    \item \textbf{Slider interactif} à 3 slides avec navigation par flèches et indicateurs à points~: présentation du projet, modélisation des alertes, et automatisation RPA
    \item \textbf{Grille de 6 fonctionnalités}~: Gestion des Réclamations, Détection d'Anomalies, Tableau de Bord, Moteur de Recommandations, Sécurité Bancaire, Intégration RPA
    \item \textbf{Section processus} illustrant le flux en 4 étapes~: Soumission Réclamation $\rightarrow$ Analyse NLP $\rightarrow$ Détection d'Anomalie $\rightarrow$ Résolution / Déclenchement Alerte RPA
    \item \textbf{Section engagement} avec 6 cartes de proposition de valeur (Intelligence, Sécurité, Disponibilité, Rapidité, Flexibilité, Simplicité)
    \item \textbf{Footer} avec branding Attijari bank et liens de navigation
    \item \textbf{Design responsive}~: support mobile avec menu hamburger dédié
\end{itemize}

\subsection{Interface Chatbot}

Le chatbot constitue le composant d'assistance et de soumission de réclamation, développé sous forme d'un \textbf{composant Angular interactif} intégrant~:

\begin{itemize}[noitemsep]
    \item \textbf{Moteur de détection d'intention NLP}~: classification par correspondance de patterns et analyse sémantique backend pour qualifier si le message est une réclamation, extraire le type d'opération (Support, Exploitation, Décisionnel...) et calculer un score d'anomalie
    \item \textbf{Architecture raccordée à l'API backend}~: intégration directe avec l'API réelle FastAPI (\texttt{http://localhost:8000/chat}), gérant de manière transparente les en-têtes JWT et les sessions
    \item \textbf{Gestion multi-sessions}~: création, suppression, persistance en \texttt{localStorage} isolée par utilisateur, renommage automatique des sessions basé sur le premier message
    \item \textbf{Panneau de contexte client} (panneau droit)~: affichage en temps réel du profil utilisateur authentifié, intention détectée, analyse sémantique, barre de confiance colorée
    \item \textbf{Sidebar latérale}~: sessions actives, boutons de requêtes rapides (consulter solde courant/épargne, bloquer sa carte bancaire, lister les transactions récentes)
    \item \textbf{Réponses enrichies}~: templates avec placeholders dynamiques pour les données bancaires réelles de l'utilisateur connecté (soldes, numéro de carte, historique de transactions)
    \item \textbf{Suivi de ticket}~: génération d'un numéro de ticket de suivi unique persistant dans PostgreSQL pour chaque réclamation soumise
    \item \textbf{Indicateur de frappe animé} avec délai aléatoire réaliste
    \item \textbf{Panneaux rétractables}~: sidebar gauche et panneau contexte droit avec animation fluide
\end{itemize}

\subsection{Tableau de Bord Opérationnel}

Le dashboard analytique est un composant complexe développé \textbf{entièrement sous forme de graphiques SVG et CSS dynamiques dans Angular}, intégrés à l'API en temps réel~:

\begin{itemize}[noitemsep]
    \item \textbf{4 cartes KPI dynamiques}~: Alertes Actives (score seuil > 0.75), Score de Risque Max LSTM, Total Tickets, et Précision de l'IA (AUC-ROC de 87\%)
    \item \textbf{Graphique en courbes} (Line/Area)~: Courbe de tendance temporelle sur 7 jours avec gradient SVG fluide pour le suivi des anomalies détectées
    \item \textbf{Graphique en barres verticales}~: Risque par Division IT, montrant la distribution des tickets par groupe d'exploitation (Support, Exploitation, Décisionnel, etc.)
    \item \textbf{Graphique Donut}~: Sévérité des Alertes, segmentant dynamiquement les alertes en 4 niveaux (Critique, Haute, Moyenne, Faible) avec le nombre d'alertes critiques au centre
    \item \textbf{Boîte de Réception \& Messagerie interactive} (inbox)~: interface permettant au responsable IT de lister les réclamations, de les filtrer, d'inspecter les détails (identité de l'utilisateur, groupe IT, score d'anomalie), de changer le statut ("Ouvert" / "Résolu") et de répondre en direct à l'utilisateur
    \item \textbf{Flux en temps réel} (Live Feed)~: événements de surveillance système rafraîchis en continu
    \item \textbf{Audit Trail complet}~: journalisation sécurisée de toutes les actions administratives et de sécurité pour une traçabilité totale
    \item \textbf{Filtres temporels et rechargement dynamique} avec configuration adaptative au redimensionnement
\end{itemize}

\subsection{Sécurité, Authentification et Contrôle d'Accès}

La sécurité et le contrôle d'accès de la solution comprennent~:

\begin{itemize}[noitemsep]
    \item \textbf{Formulaire d'authentification réactif}~: Saisie sécurisée des identifiants, communication asynchrone avec l'API FastAPI et stockage sécurisé du jeton JWT dans le local storage.
    \item \textbf{Formulaire d'inscription complet}~: Validation en temps réel des champs (CIN unique, format email, force du mot de passe).
    \item \textbf{Garde d'accès Angular} (\texttt{auth.guard.ts})~: Intercepte les accès aux routes privées (Dashboard, Alertes) et vérifie la validité du jeton JWT via le rôle (RBAC) de l'utilisateur.
    \item \textbf{Service d'authentification centralisé} (\texttt{AuthService})~: Gère l'état global de l'utilisateur connecté via un \texttt{BehaviorSubject} Angular, synchronisant l'état visuel en temps réel.
\end{itemize}

\section{API REST Backend (FastAPI)}

L'API REST backend a été développée sous le framework asynchrone FastAPI (Python 3.11) pour gérer l'ensemble des opérations métiers, de l'authentification et de la détection d'anomalies. Elle expose 8 endpoints REST sécurisés raccordés à la base PostgreSQL et aux modèles IA~:

\begin{table}[H]
\centering
\caption{Endpoints de l'API REST Réelle (FastAPI)}
\begin{tabularx}{\textwidth}{|l|l|X|}
\hline
\textbf{Endpoint} & \textbf{Méthode} & \textbf{Fonctionnalité} \\
\hline
\texttt{/api/auth/register} & POST & Inscription d'un utilisateur, validation des champs, hachage bcrypt \\
\hline
\texttt{/api/auth/login} & POST & Authentification, génération du token JWT sécurisé \\
\hline
\texttt{/api/auth/me} & GET & Récupération du profil de l'utilisateur connecté via JWT \\
\hline
\texttt{/api/reclamations} & GET/POST & Liste de toutes les réclamations / Soumission d'une réclamation NLP \\
\hline
\texttt{/api/reclamations/\{id\}/status} & PATCH & Mise à jour du statut d'un ticket (Ouvert / Résolu) par le Responsable IT \\
\hline
\texttt{/api/reclamations/\{id\}/reply} & POST & Envoi d'une réponse de l'admin et notification par email automatique \\
\hline
\texttt{/api/predictions/lstm} & GET & Analyse prédictive des anomalies temporelles via le modèle LSTM \\
\hline
\texttt{/api/recommandations/knn} & GET & Suggestions d'actions correctives basées sur la similarité KNN \\
\hline
\end{tabularx}
\end{table}

\section{Base de Données PostgreSQL}

Le schéma relationnel est hébergé sous PostgreSQL pour garantir la cohérence des données, l'intégrité référentielle et de hautes performances~:

\begin{itemize}[noitemsep]
    \item \textbf{Table \texttt{users}}~: Stocke les informations des utilisateurs (CIN unique, nom, email, mot de passe haché par bcrypt, rôle admin/user).
    \item \textbf{Table \texttt{reclamations}}~: Enregistre l'ensemble des tickets de réclamations (id unique, date de création, objet, description sémantique, division IT détectée, score d'anomalie et statut).
    \item \textbf{Table \texttt{audit\_logs}}~: Enregistre l'ensemble des actions administratives et de sécurité pour une traçabilité complète.
    \item \textbf{Persistance robuste}~: Base PostgreSQL hébergée en local avec gestion de sessions asynchrones via SQLAlchemy et psycopg2.
\end{itemize}

\section{Architecture Globale de la Solution}

L'architecture globale de la solution assure un découplage complet et une communication fluide entre les différentes couches de l'application~:

\subsection{Frontend Angular 21 (Single Page Application)}

La couche de présentation offre une interface réactive et modulaire organisée autour des éléments suivants~:

\begin{itemize}[noitemsep]
    \item \textbf{Composants de page}~: \texttt{HomeComponent} (accueil), \texttt{LoginComponent} (connexion), \texttt{RegisterComponent} (inscription), \texttt{DashboardComponent} (tableau de bord de détection d'anomalies), \texttt{ChatComponent} (chatbot interactif).
    \item \textbf{Composants transversaux}~: Barre de navigation supérieure (\texttt{NavbarComponent}), pied de page (\texttt{FooterComponent}), boîte de réclamations interactive (\texttt{Inbox}).
    \item \textbf{Services Angular injectables}~: \texttt{AuthService} (sécurité JWT), \texttt{ChatService} (communication chatbot) et \texttt{TicketService} (gestion des réclamations).
    \item \textbf{Système de Routage}~: Routage réactif Angular (\texttt{app.routes.ts}) avec gardes d'accès (guards) empêchant les accès non autorisés au dashboard.
\end{itemize}

\subsection{Backend FastAPI (Python 3.11)}

La couche de services asynchrones sous FastAPI assure le traitement métier et l'intégration des modèles IA~:

\begin{itemize}[noitemsep]
    \item \textbf{FastAPI Web Framework}~: Gestion asynchrone ultra-rapide des requêtes REST (Uvicorn).
    \item \textbf{SQLAlchemy ORM}~: Couche d'accès aux données PostgreSQL avec sessions asynchrones.
    \item \textbf{Sécurité JWT}~: Hachage bcrypt des mots de passe et validation des en-têtes JWT.
    \item \textbf{psycopg2 Driver}~: Pilote de base de données PostgreSQL natif haute performance.
    \item \textbf{Organisation du code}~: Architecture modulaire propre en modules (\texttt{app.routers}, \texttt{app.models}, \texttt{app.database}, \texttt{app.services}).
\end{itemize}

\section{Design System}

Un design system complet a été élaboré, intégrant~:

\begin{itemize}[noitemsep]
    \item \textbf{Palette Attijari}~: Orange primaire (\texttt{\#F7941D}), fond sombre (\texttt{\#141416}), textes hiérarchisés
    \item \textbf{Typographie}~: Inter de Google Fonts, 7 graisses (300--900)
    \item \textbf{Effets visuels}~: glassmorphism (\texttt{backdrop-filter: blur}), gradients, ombres multicouches
    \item \textbf{Micro-animations}~: transitions hover, compteurs animés, indicateurs pulsants
    \item \textbf{Responsivité}~: breakpoints mobile, tablette, desktop
    \item \textbf{Variables CSS}~: système de tokens pour les couleurs, espacements, transitions
\end{itemize}


% ============================================================
% CHAPTER 3: TRAVAUX RÉALISÉS — AXE IA & DONNÉES
% ============================================================
\chapter{Travaux Réalisés --- Axe Intelligence Artificielle \& Données}

L'axe Intelligence Artificielle constitue le c\oe ur analytique de la solution. Il s'agit de transformer les données brutes de réclamations bancaires en intelligence opérationnelle exploitable.

\section{Réception et Exploration des Données}

Le fichier CSV de réclamations fourni par la banque a été réceptionné en début de projet. Une exploration initiale a permis d'identifier la structure du dataset, les colonnes disponibles, les types de données, et les premières statistiques descriptives.

\section{Traitement et Nettoyage des Données}

Le pipeline de nettoyage, développé en Python avec les bibliothèques \texttt{pandas} et \texttt{numpy}, a été appliqué selon les étapes suivantes~:

\begin{table}[H]
\centering
\caption{Pipeline de nettoyage des données}
\begin{tabularx}{\textwidth}{|c|X|l|}
\hline
\textbf{Étape} & \textbf{Opération} & \textbf{Outil} \\
\hline
1 & Suppression des enregistrements dupliqués & pandas \\
\hline
2 & Traitement des valeurs manquantes (NaN) par colonne & pandas \\
\hline
3 & Correction des types de données (dates, numériques, catégories) & pandas \\
\hline
4 & Normalisation des descriptions textuelles (casse, accents, ponctuation) & regex \\
\hline
5 & Détection des valeurs aberrantes (méthode IQR) & numpy · scipy \\
\hline
6 & Analyse exploratoire (EDA) avec graphiques & matplotlib · seaborn \\
\hline
\end{tabularx}
\end{table}

\noindent Les résultats du nettoyage sont~:

\begin{table}[H]
\centering
\caption{Résultats du nettoyage}
\begin{tabularx}{\textwidth}{|X|l|l|}
\hline
\textbf{Métrique} & \textbf{Avant} & \textbf{Après} \\
\hline
Enregistrements dupliqués & Présents & Supprimés \\
\hline
Valeurs manquantes & Multiples colonnes & Traitées \\
\hline
Types incohérents & Dates en string & Corrigés \\
\hline
Valeurs aberrantes & Non traitées & Détectées et traitées (IQR) \\
\hline
Descriptions textuelles & Casse mixte & Normalisées \\
\hline
\end{tabularx}
\end{table}

Le dataset nettoyé a été stocké dans la base de données, prêt pour l'indexation et l'entraînement des modèles.

\section{Pipeline NLP (En cours)}

Le traitement du langage naturel permet au système de comprendre le contenu sémantique de chaque réclamation bancaire~:

\begin{itemize}[noitemsep]
    \item \textbf{Pipeline spaCy}~: tokenisation, lemmatisation, extraction d'entités nommées (NER) pour identifier les types d'opérations, montants, dates et références
    \item \textbf{Embeddings BERT}~: génération de vecteurs de 768 dimensions par réclamation via \texttt{sentence-transformers}, capturant le sens profond du texte
    \item \textbf{Score d'anomalie}~: calcul par similarité cosinus entre chaque nouvelle réclamation et l'historique, permettant d'identifier les cas atypiques
    \item \textbf{Indexation Elasticsearch}~: recherche full-text sur le corpus de réclamations pour l'interface de filtre et de recherche
\end{itemize}

Le moteur de détection d'intention implémenté côté frontend sert de \textbf{couche de transition} vers le modèle NLP complet du backend. L'architecture du chatbot permet de basculer du mode démonstration vers l'API réelle en modifiant un paramètre de configuration unique.

\section{Modèle Prédictif LSTM (En cours)}

Le modèle de prédiction LSTM (Long Short-Term Memory) constitue la valeur ajoutée majeure du projet~:

\begin{itemize}[noitemsep]
    \item \textbf{Architecture}~: réseau LSTM entraîné sur des séquences temporelles de réclamations glissant sur 30 jours
    \item \textbf{Sortie}~: score de risque entre 0 et 1 pour les 24 heures à venir, par type de processus bancaire
    \item \textbf{Seuil d'alerte}~: configurable (défaut~: 0.75) --- déclenche une notification si dépassé
    \item \textbf{Versioning}~: chaque modèle entraîné sera versionné dans MLflow avec ses métriques
    \item \textbf{Framework}~: TensorFlow avec optimisation par grille des hyperparamètres
\end{itemize}

\section{Sécurité des Données}

La sécurité des données bancaires est assurée par plusieurs mécanismes~:

\begin{table}[H]
\centering
\caption{Mécanismes de sécurité implémentés}
\begin{tabularx}{\textwidth}{|l|X|c|}
\hline
\textbf{Mécanisme} & \textbf{Implémentation} & \textbf{Statut} \\
\hline
Hachage mots de passe & bcrypt (coût 12) & \done \\
\hline
Authentification & Sessions sécurisées, migration JWT prévue & \done \\
\hline
Contrôle d'accès & RBAC~: admin, utilisateur & \done \\
\hline
Chiffrement en base & AES-256 pour données sensibles & \done \\
\hline
Audit trail & Journalisation des actions & \inprog \\
\hline
Sauvegarde & Backup quotidien automatique (prévu) & \inprog \\
\hline
\end{tabularx}
\end{table}


% ============================================================
% CHAPTER 4: OBJECTIFS
% ============================================================
\chapter{Objectifs}

\section{Objectifs Globaux du Projet}

\begin{table}[H]
\centering
\caption{Objectifs globaux}
\begin{tabularx}{\textwidth}{|c|l|X|}
\hline
\textbf{\#} & \textbf{Objectif} & \textbf{Description} \\
\hline
O1 & Analyser les réclamations & Extraction d'informations par NLP/BERT sur le texte \\
\hline
O2 & Prédire les incidents & Modèle LSTM anticipant les incidents 24h à l'avance \\
\hline
O3 & Recommander les actions & Moteur KNN suggérant les solutions basées sur l'historique \\
\hline
O4 & Automatiser l'exécution & Robots RPA (UiPath) pour les corrections P1-P2 \\
\hline
O5 & Interface web sécurisée & Tableau de bord interactif pour les équipes opérationnelles \\
\hline
O6 & Apprentissage continu & Réentraînement hebdomadaire avec les nouveaux cas \\
\hline
\end{tabularx}
\end{table}

\section{Bilan de Mi-Parcours (Semaines 1--4)}

\begin{table}[H]
\centering
\caption{État d'avancement à mi-parcours}
\begin{tabularx}{\textwidth}{|X|c|l|}
\hline
\textbf{Objectif mi-parcours} & \textbf{Statut} & \textbf{Avancement} \\
\hline
Interface frontend complète (6 pages) & \done & 100\% \\
\hline
Base de données opérationnelle & \done & 100\% \\
\hline
API REST fonctionnelle (6 endpoints) & \done & 100\% \\
\hline
Design system complet (responsive, animations) & \done & 100\% \\
\hline
Système d'authentification RBAC & \done & 100\% \\
\hline
Nettoyage et stockage des données bancaires & \done & 100\% \\
\hline
Migration Angular amorcée & \inprog & 70\% \\
\hline
Pipeline NLP (spaCy + BERT) & \inprog & En cours \\
\hline
Modèle prédictif LSTM & \inprog & En cours \\
\hline
Backend FastAPI opérationnel & \done & 100\% \\
\hline
\end{tabularx}
\end{table}

\section{Objectifs des Semaines 5--8}

\begin{table}[H]
\centering
\caption{Objectifs pour la seconde moitié du projet}
\begin{tabularx}{\textwidth}{|X|l|c|}
\hline
\textbf{Objectif} & \textbf{Semaines} & \textbf{Priorité} \\
\hline
Finaliser le pipeline NLP et le modèle LSTM & S5 & Haute \\
\hline
Développer le moteur de recommandations KNN & S5 & Haute \\
\hline
Intégrer UiPath Orchestrator pour l'automatisation RPA & S5--S6 & Haute \\
\hline
Développer l'API FastAPI complète avec WebSocket & S5--S6 & Haute \\
\hline
Connecter le frontend aux données réelles de l'API IA & S5--S6 & Haute \\
\hline
Déploiement et tests de l'API FastAPI & S5--S6 & Haute \\
\hline
Tests unitaires et tests de sécurité & S7 & Haute \\
\hline
Rédaction du rapport final et préparation soutenance & S7--S8 & Haute \\
\hline
\end{tabularx}
\end{table}


% ============================================================
% CHAPTER 5: MÉTHODOLOGIE
% ============================================================
\chapter{Méthodologie}

\section{Approche de Développement}

Le projet adopte une \textbf{méthodologie agile adaptée} au contexte académique, organisée en 4 phases itératives sur 8 semaines. Chaque phase produit un livrable concret validable par l'encadrant et se termine par un jalon de contrôle.

\begin{table}[H]
\centering
\caption{Phases du projet}
\begin{tabularx}{\textwidth}{|c|l|X|l|}
\hline
\textbf{Phase} & \textbf{Période} & \textbf{Activités} & \textbf{Livrable} \\
\hline
1 & 23 Mars -- 5 Avril & Interface frontend conforme Attijari · Création BDD · Nettoyage données · Stockage & Interface + BDD \\
\hline
2 & 6 Avril -- 19 Avril & Python NLP + Prédictions · Frontend avancé & Pipeline NLP + LSTM \\
\hline
3 & 20 Avril -- 3 Mai & UiPath (robots RPA) · Frontend + API FastAPI & RPA + API \\
\hline
4 & 4 Mai -- 17 Mai & Finalisation · Tests · Rapport · Soutenance & Projet complet \\
\hline
\end{tabularx}
\end{table}

\vspace{0.3cm}
\begin{center}
\begin{tikzpicture}[x=0.85cm]
    % Timeline
    \draw[thick, ->, attijariDark] (0,0) -- (17.5,0);
    
    % Phase boxes
    \fill[attijariOrange!20] (0,-0.4) rectangle (5.2,0.4);
    \fill[attijariOrange!35] (5.4,-0.4) rectangle (10.2,0.4);
    \fill[attijariOrange!15] (10.4,-0.4) rectangle (14.2,0.4);
    \fill[attijariOrange!10] (14.4,-0.4) rectangle (17.2,0.4);
    
    % Phase labels
    \node[font=\scriptsize\bfseries] at (2.6,0) {Phase 1};
    \node[font=\scriptsize\bfseries] at (7.8,0) {Phase 2};
    \node[font=\scriptsize\bfseries] at (12.3,0) {Phase 3};
    \node[font=\scriptsize\bfseries] at (15.8,0) {Phase 4};
    
    % Dates
    \node[font=\tiny, below] at (0,-0.5) {23/03};
    \node[font=\tiny, below] at (5.3,-0.5) {05/04};
    \node[font=\tiny, below] at (10.3,-0.5) {19/04};
    \node[font=\tiny, below] at (14.3,-0.5) {03/05};
    \node[font=\tiny, below] at (17.2,-0.5) {17/05};
    
    % Mi-parcours marker
    \draw[thick, red, dashed] (8.5,-1) -- (8.5,1);
    \node[font=\tiny\bfseries, red, above] at (8.5,1) {MI-PARCOURS};
    \node[font=\tiny, red, below] at (8.5,-1) {17/04};
\end{tikzpicture}
\end{center}

\section{Tâches Détaillées par Phase}

\subsection{Phase 1 --- Interface et Données (23 Mars -- 5 Avril) \done}

\begin{enumerate}[noitemsep]
    \item Analyser le sujet 21 en détail et identifier les livrables attendus
    \item Rencontrer l'encadrant bancaire, clarifier les attentes et confirmer le périmètre
    \item Concevoir les maquettes d'interface conformes à la charte graphique Attijari bank
    \item Développer la page d'accueil avec hero, slider, fonctionnalités et processus
    \item Développer l'interface chatbot avec détection d'intention et gestion multi-sessions
    \item Développer le tableau de bord avec graphiques SVG natifs sous Angular 21 et flux en temps réel
    \item Développer les pages de connexion et d'inscription avec validation
    \item Implémenter le module d'authentification RBAC partagé (guards, interceptors)
    \item Concevoir le schéma de la base de données PostgreSQL (tables, relations et contraintes FK)
    \item Développer l'API REST FastAPI avec authentification bcrypt et jetons JWT
    \item Réceptionner et explorer le fichier CSV de réclamations bancaires
    \item Appliquer le pipeline de nettoyage (doublons, NaN, outliers, normalisation)
    \item Stocker le dataset nettoyé en base de données
    \item Élaborer le design system complet (palette, typographie, animations, responsive)
\end{enumerate}

\subsection{Phase 2 --- NLP, Prédictions et Frontend avancé (6 -- 19 Avril) \inprog}

\begin{enumerate}[noitemsep, start=15]
    \item Construire le pipeline NLP avec spaCy (tokenisation, NER, lemmatisation)
    \item Générer les embeddings BERT (768 dimensions) pour chaque réclamation
    \item Calculer les scores d'anomalie par similarité cosinus
    \item Indexer les réclamations dans Elasticsearch
    \item Préparer les séquences temporelles (fenêtre 30 jours) pour le LSTM
    \item Entraîner le modèle LSTM avec TensorFlow
    \item Amorcer la migration du frontend vers Angular 21 (5 pages, routing, services)
    \item Développer l'architecture de base du backend FastAPI
\end{enumerate}

\subsection{Phase 3 --- UiPath, Frontend et API (20 Avril -- 3 Mai)}

\begin{enumerate}[noitemsep, start=23]
    \item Configurer UiPath Orchestrator et créer les robots pour les actions P1-P2
    \item Développer l'API FastAPI complète avec authentification JWT
    \item Configurer les WebSockets pour les alertes temps réel
    \item Documenter l'API via Swagger UI
    \item Connecter le frontend Angular aux données réelles de l'API IA
    \item Développer le moteur de recommandations KNN
    \item Réaliser les tests d'intégration bout en bout
\end{enumerate}

\subsection{Phase 4 --- Finalisation (4 -- 17 Mai)}

\begin{enumerate}[noitemsep, start=30]
    \item Écrire les tests unitaires (couverture $>$ 70\%)
    \item Tests de sécurité (injection SQL, XSS, validation JWT)
    \item Rédiger la partie implémentation du rapport avec captures d'écran
    \item Rédiger la conclusion et les perspectives d'évolution
    \item Préparer et répéter la présentation de soutenance
    \item Démonstration live du système complet
\end{enumerate}

\section{Outils et Environnement de Développement}

\begin{table}[H]
\centering
\caption{Environnement technique}
\begin{tabularx}{\textwidth}{|l|X|}
\hline
\textbf{Catégorie} & \textbf{Outil} \\
\hline
Éditeur de code & Visual Studio Code (extensions Angular, TypeScript, Python, GitLens) \\
\hline
Gestion de versions & Git avec GitHub (dépôt privé, branches fonctionnelles) \\
\hline
Serveur local & Uvicorn (FastAPI Backend) · Angular Dev Server (Frontend) \\
\hline
Base de données & PostgreSQL · pgAdmin 4 \\
\hline
Build frontend & Angular CLI 21 · npm · TypeScript 5.9 \\
\hline
Serveur Backend & Python 3.11 · FastAPI · Uvicorn \\
\hline
Backend IA & Python · FastAPI · TensorFlow · spaCy · scikit-learn \\
\hline
Conteneurisation & Docker Desktop (Elasticsearch, Redis, MLflow) \\
\hline
Automatisation & UiPath Studio · UiPath Orchestrator \\
\hline
Documentation API & Swagger UI (auto-généré par FastAPI) \\
\hline
\end{tabularx}
\end{table}


% ============================================================
% CHAPTER 6-7: RÉSULTATS
% ============================================================
\chapter{Résultats --- Axe Développement}

\section{Interface Utilisateur Livrée}

L'ensemble de l'interface web de la solution est fonctionnel et opérationnel~:

\begin{table}[H]
\centering
\caption{Bilan quantitatif du développement frontend}
\begin{tabularx}{\textwidth}{|l|c|X|}
\hline
\textbf{Composant} & \textbf{Volume} & \textbf{Fonctionnalités clés} \\
\hline
\texttt{index.html} & 454 lignes & Hero animé, slider, 6 features, processus, footer \\
\hline
\texttt{chat.html} + \texttt{chat.ts} & 540 lignes & Composant chatbot Angular, routage, assistance, panneau contexte \\
\hline
\texttt{dashboard.html} + \texttt{dashboard.ts} & 940 lignes & 4 KPI, tendances SVG, barres divisions IT, donut sévérité, boîte de réclamations interactive (inbox) \\
\hline
\texttt{login.html} + \texttt{login.ts} & 220 lignes & Formulaire d'authentification Angular, sécurité JWT, redirection \\
\hline
\texttt{register.html} + \texttt{register.ts} & 340 lignes & Sections d'inscription, validation de formulaires réactifs \\
\hline
\texttt{auth.guard.ts} & 80 lignes & Contrôle d'accès et sécurité RBAC sur les routes Angular \\
\hline
\texttt{app.css} + \texttt{dashboard.css} & $\sim$1200 lignes & Design system, charte graphique sombre, animations, responsive \\
\hline
API REST (FastAPI) & 8 fichiers & FastAPI, SQLAlchemy ORM, endpoints authentification et tickets \\
\hline
Base de données & PostgreSQL & Schéma PostgreSQL, tables relationnelles, contraintes FK \\
\hline
\textbf{Total V1} & \textbf{$\sim$4\,600} & \textbf{Solution full-stack fonctionnelle} \\
\hline
\end{tabularx}
\end{table}

\section{Migration Architecturale}

\begin{table}[H]
\centering
\caption{État de la migration Angular}
\begin{tabularx}{\textwidth}{|l|c|X|}
\hline
\textbf{Composant} & \textbf{Statut} & \textbf{Détails} \\
\hline
\texttt{HomeComponent} & \done & Hero + features fidèles aux spécifications \\
\hline
\texttt{ChatComponent} & \done & Chatbot intégré en Angular \\
\hline
\texttt{DashboardComponent} & \done & Canvas charts migrés \\
\hline
\texttt{LoginComponent} & \done & Formulaire réactif \\
\hline
\texttt{RegisterComponent} & \done & Multi-sections \\
\hline
\texttt{NavbarComponent} & \done & Navigation SPA routerLink \\
\hline
\texttt{FooterComponent} & \done & Réutilisable \\
\hline
\texttt{AuthService} & \done & Gestion d'état, HTTPClient \\
\hline
\texttt{ChatService} & \done & Communication API \\
\hline
Routing & \done & 5 routes + wildcard \\
\hline
\end{tabularx}
\end{table}

\section{Indicateurs de Qualité}

\begin{table}[H]
\centering
\caption{Métriques de qualité}
\begin{tabularx}{\textwidth}{|X|l|}
\hline
\textbf{Métrique} & \textbf{Valeur} \\
\hline
Pages développées sous Angular 21 & 5 \\
\hline
Composants Angular & 7 (5 pages + 2 réutilisables) \\
\hline
Services Angular & 2 (Auth + Chat) \\
\hline
Endpoints API & 6 \\
\hline
Charts Canvas natifs & 4 (line, bar, donut, stacked) \\
\hline
Compatibilité navigateurs & Chrome, Firefox, Edge \\
\hline
Design responsive & Mobile, tablette, desktop \\
\hline
\end{tabularx}
\end{table}


\chapter{Résultats --- Axe Intelligence Artificielle}

\section{Nettoyage des Données}

Le pipeline de nettoyage a été exécuté avec succès sur le dataset CSV de réclamations bancaires fourni par Attijari bank. Toutes les étapes de suppression des doublons, traitement des valeurs manquantes, correction des types, normalisation des textes et détection des outliers par la méthode IQR ont été complétées.

Le dataset nettoyé a été importé et stocké dans la base de données, constituant la fondation pour l'entraînement des modèles d'intelligence artificielle.

\section{Pipeline NLP}

\begin{table}[H]
\centering
\caption{État du pipeline NLP}
\begin{tabularx}{\textwidth}{|l|X|c|}
\hline
\textbf{Composant} & \textbf{Résultat} & \textbf{Statut} \\
\hline
Pipeline spaCy & Tokenisation, lemmatisation, NER opérationnels & \inprog \\
\hline
Embeddings BERT & Vecteurs 768 dimensions par réclamation & \inprog \\
\hline
Scores d'anomalie & Similarité cosinus avec historique & \inprog \\
\hline
Elasticsearch & Indexation pour recherche full-text & \inprog \\
\hline
\end{tabularx}
\end{table}

\section{Modèle Prédictif LSTM}

\begin{table}[H]
\centering
\caption{Avancement du modèle LSTM}
\begin{tabularx}{\textwidth}{|X|c|}
\hline
\textbf{Aspect} & \textbf{État} \\
\hline
Préparation des séquences temporelles (fenêtre 30 jours) & \inprog \\
\hline
Architecture LSTM définie & \inprog \\
\hline
Entraînement initial avec TensorFlow & \inprog \\
\hline
Versioning MLflow (prévu) & Planifié \\
\hline
Score de risque 0$\rightarrow$1 par processus & \inprog \\
\hline
\end{tabularx}
\end{table}


% ============================================================
% CHAPTER 8: DIFFICULTÉS
% ============================================================
\chapter{Difficultés Rencontrées}

\section{Difficultés Techniques}

\begin{table}[H]
\centering
\caption{Difficultés techniques et solutions}
\begin{tabularx}{\textwidth}{|l|X|X|}
\hline
\textbf{Difficulté} & \textbf{Impact} & \textbf{Solution adoptée} \\
\hline
Graphiques SVG dynamiques réactifs & Rendre les graphiques fluides et adaptables au redimensionnement sans bibliothèque tierce & Modélisation géométrique des courbes et barres en éléments SVG réactifs liés à l'état Angular \\
\hline
Synchronisation authentification & Le chatbot doit connaître l'utilisateur pour personnaliser les réponses & Utilisation d'un BehaviorSubject centralisé et injection du profil utilisateur connecté \\
\hline
Intégration multi-technologies & Raccorder le frontend Angular à l'API FastAPI en temps réel & Architecture modulaire en couches, CORS configuré avec sécurité JWT \\
\hline
Données sensibles & Impossibilité d'utiliser des données client réelles & Données synthétiques réalistes \\
\hline
Qualité du CSV & Doublons, NaN, outliers, encodage incohérent & Pipeline en 6 étapes \\
\hline
\end{tabularx}
\end{table}

\section{Difficultés Organisationnelles}

\begin{table}[H]
\centering
\caption{Difficultés organisationnelles et solutions}
\begin{tabularx}{\textwidth}{|l|X|X|}
\hline
\textbf{Difficulté} & \textbf{Impact} & \textbf{Solution adoptée} \\
\hline
Ampleur du périmètre & Le projet couvre frontend, backend, BDD, NLP, LSTM, KNN et RPA sur une seule personne & Planning serré avec priorisation stricte et itérations rapides \\
\hline
Disponibilité tardive des données & CSV fourni tardivement & Avancement du frontend en parallèle \\
\hline
Complexité technologique variée & Maîtrise simultanée de TypeScript, Python, SQL, et RPA & Documentation abondante et apprentissage continu \\
\hline
Coordination avec la banque & Contraintes de confidentialité et d'accès à l'intranet & Working sessions régulières avec l'encadrant \\
\hline
\end{tabularx}
\end{table}


% ============================================================
% CHAPTER 9: PERSPECTIVES
% ============================================================
\chapter{Perspectives}

\section{Perspectives à Court Terme (Semaines 5--8)}

\begin{table}[H]
\centering
\caption{Perspectives immédiates}
\begin{tabularx}{\textwidth}{|X|l|}
\hline
\textbf{Perspective} & \textbf{Période} \\
\hline
Finalisation du modèle LSTM --- optimisation hyperparamètres, validation croisée, déploiement du meilleur modèle via MLflow & S5 \\
\hline
Développement du moteur de recommandations KNN --- calcul du taux de succès par action corrective & S5 \\
\hline
Intégration de l'automatisation RPA via UiPath Orchestrator --- robots pour les actions P1-P2 & S5--S6 \\
\hline
Développement de l'API FastAPI complète avec WebSocket pour les alertes temps réel & S5--S6 \\
\hline
Connexion du frontend aux données réelles --- basculement du mode démonstration vers l'API & S5--S6 \\
\hline
Enrichissement du dashboard avec les données LSTM réelles (score de risque, alertes, recommandations) & S6--S7 \\
\hline
Tests unitaires (couverture $>$ 70\%), tests de sécurité et d'intégration & S7 \\
\hline
Rédaction du rapport final et préparation de la soutenance & S7--S8 \\
\hline
\end{tabularx}
\end{table}

\section{Perspectives d'Évolution Post-PFE}

\begin{itemize}[noitemsep]
    \item \textbf{Application mobile}~: développement d'une application native (React Native / Flutter) pour l'accès client mobile
    \item \textbf{NLP multilingue}~: extension du pipeline pour supporter l'arabe tunisien (dialecte) et l'arabe standard, en plus du français
    \item \textbf{Intégration vocale}~: module de reconnaissance vocale pour le call center téléphonique
    \item \textbf{Déploiement cloud}~: migration vers une infrastructure cloud (AWS / Azure) pour la scalabilité
    \item \textbf{Tableau de bord prédictif avancé}~: carte de chaleur des processus à risque et module de simulation what-if
    \item \textbf{Boucle d'apprentissage avancée}~: reinforcement learning pour optimiser automatiquement les seuils d'alerte
    \item \textbf{Intégration omnicanale}~: extension vers les canaux SMS, WhatsApp et email pour une couverture complète
\end{itemize}


% ============================================================
% CHAPTER 10: CONCLUSION
% ============================================================
\chapter{Conclusion}

Ce rapport de mi-parcours rend compte de l'avancement significatif réalisé au cours des quatre premières semaines du projet de fin d'études au sein d'Attijari bank.

Sur l'axe développement, un \textbf{système frontend complet et fonctionnel} a été livré sous Angular, comprenant le chatbot d'assistance et de soumission de réclamation, un tableau de bord analytique de détection d'anomalies avec graphiques SVG dynamiques (tendances 7 jours, divisions IT, donut sévérité), une boîte de réclamations interactive (inbox) avec messagerie en direct, un système d'authentification RBAC et une API REST opérationnelle sous FastAPI et PostgreSQL.

Sur l'axe intelligence artificielle, les fondations sont solidement posées~: le pipeline de nettoyage des données est achevé, le dataset nettoyé est stocké en base, et le pipeline NLP (spaCy + BERT) ainsi que le modèle prédictif LSTM sont en cours de développement et d'optimisation.

Les principales difficultés rencontrées --- complexité des graphiques dynamiques sous Angular, intégration en temps réel des flux de surveillance, et ampleur du périmètre technique --- ont été surmontées par des solutions techniques appropriées et une organisation rigoureuse du travail.

Pour la seconde moitié du projet (semaines 5 à 8), les priorités sont clairement définies~: finalisation du modèle LSTM et du moteur de recommandations KNN, intégration de l'automatisation RPA via UiPath, développement de l'API FastAPI complète, connexion du frontend aux données réelles, et préparation du livrable final pour la soutenance.

Le projet est dans l'ensemble \textbf{conforme au planning prévisionnel} et les objectifs de mi-parcours ont été atteints. La seconde phase s'annonce intensive mais réalisable dans les délais impartis, avec pour objectif la livraison d'un système complet et fonctionnel le 17 mai 2026.

\vspace{1cm}
\begin{center}
\rule{0.5\textwidth}{0.4pt}\\[0.3cm]
{\small\textcolor{attijariGray}{Document rédigé dans le cadre du PFE 2026}}\\
{\small\textcolor{attijariGray}{Attijari bank --- Sujet 21}}\\
{\small\textcolor{attijariGray}{Spécialité Génie Logiciel}}\\
{\small\textcolor{attijariGray}{Rationalisation, Reengineering \& Efficacité Opérationnelle}}
\end{center}

\end{document}
